\section{Finite syntax and proofs}
\label{app:fork-proofs}

Let $\mathcal E=\mathcal C_{\mathrm{rec}}^{J}$, where
$J=(1\leftarrow0\to2)$, and let
$\mathcal A=\operatorname{Par}(\mathcal E,\operatorname{Mono})$.
The components of the atom object $V$ are $V_0=\N^4$ and
$V_1=V_2=\N^2$; write $r_i=r_i^V=V(0\to i):V_0\to V_i$
for $i=1,2$.
We construct formula and proof objects in $\mathcal E$. The image of
the conclusion map determines a restriction idempotent in $\mathcal A$.

\subsection{The syntax algebra}

Fix recursive encodings of finite rooted ordered trees and of finite
tuples of natural numbers, with recursive parsing operations. For
$j\in\{0,1,2\}$, let $S_j$ be the set of finite formula trees
generated by
\[
 \theta_j(x)\quad(x\in V_j),\qquad
 \nu_j(\varphi)\quad(\varphi\in S_j),\qquad
 \varphi\Rightarrow_j\psi\quad(\varphi,\psi\in S_j).
\]
The three constructor tags are distinct. Tree parsing and the
decidable numerical carriers $V_j$ make $S_j$ a recursive set of
codes. For $i=1,2$, define $\sigma_i:S_0\to S_i$ by replacing
each atom label $x$ with $r_i(x)$ and retaining every constructor.
These total recursive functions define the functor
$S:J\to\mathcal C_{\mathrm{rec}}$ by $S(j)=S_j$ and
$S(0\to i)=\sigma_i$. The constructors form natural transformations
$\theta:V\to S$, $\nu:S\to S$, and
$\Rightarrow:S\times S\to S$, hence morphisms of $\mathcal E$.

The covariant functor $P:\mathcal E\to\mathcal E$ is defined by
\[
 P(Z)=V+Z+(Z\times Z),\qquad
 P(f)=1_V+f+(f\times f)
 \quad(f:Z\to Z'\text{ in }\mathcal E).
\]
Its identity and composition laws follow componentwise from those
of tagged sums and ordered pairs. The constructor morphisms define
$\iota=[\theta,\nu,\Rightarrow]:P(S)\to S$.

\begin{proposition}\label{prop:fork-syntax-initial}
The algebra $\iota:P(S)\to S$ is initial in the category of
$P$-algebras in $\mathcal E$: for every morphism
$\alpha:P(Z)\to Z$ of $\mathcal E$, there is a unique morphism
$u:S\to Z$ of $\mathcal E$ satisfying $u\iota=\alpha P(u)$.
\end{proposition}
\begin{proof}
At component $j$, define $u_j:S_j\to Z_j$ by evaluating a tree
from its leaves to its root, using the composites of
$\alpha_j:V_j+Z_j+Z_j^2\to Z_j$ with the three coproduct injections.
Induction on tree size proves termination and membership
of the output in $Z_j$. The evaluation is partial recursive on
numerical codes and total on $S_j$, as required for a morphism of
$\mathcal C_{\mathrm{rec}}$.

Write $z_i=Z(0\to i):Z_0\to Z_i$ for $i=1,2$. Naturality of
$\alpha$ says that applying $z_i$ after one constructor operation
is the same as first applying $r_i$ to an atom label, or $z_i$
to each recursive argument, and then using the corresponding
operation at component $i$. Induction on the tree consequently
gives $z_i u_0=u_i\sigma_i$. Thus the component maps form
$u:S\to Z$ in $\mathcal E$, and their recursive definition gives
the algebra equation. Any morphism satisfying that equation agrees
with $u$ on atoms and then on negations and implications by the
same finite induction, proving uniqueness.
\end{proof}

\subsection{The proof object and its range}

Let $a,b$ be disjoint restriction idempotents on $V$ in $\mathcal A$.
Let $A,B\subseteq V$ be the c.e.\ subfunctors on which
$a$ and $b$ are the partial identities, with component subsets
$A_j,B_j\subseteq V_j$. Disjointness
means $A_j\cap B_j=\varnothing$ at each component.
For formulas $\alpha,\beta,\gamma\in S_j$, put
\begin{align*}
 H_{1,j}(\alpha,\beta)
   &=\alpha\Rightarrow_j(\beta\Rightarrow_j\alpha),\\
 H_{2,j}(\alpha,\beta,\gamma)
   &=(\alpha\Rightarrow_j(\beta\Rightarrow_j\gamma))
     \Rightarrow_j((\alpha\Rightarrow_j\beta)
                          \Rightarrow_j(\alpha\Rightarrow_j\gamma)),\\
 H_{3,j}(\alpha,\beta)
   &=(\nu_j\beta\Rightarrow_j\nu_j\alpha)
                          \Rightarrow_j(\alpha\Rightarrow_j\beta).
\end{align*}
A proof at component $j$ is a finite rooted ordered tree with
the following node data and conclusion. A leaf is either a
tagged logical instance $H_{1,j}(\alpha,\beta)$,
$H_{2,j}(\alpha,\beta,\gamma)$, or
$H_{3,j}(\alpha,\beta)$, carrying its displayed formula
parameters; a positive leaf labeled by $x\in A_j$, with
conclusion $\theta_j(x)$; or a negative leaf labeled by
$x\in B_j$, with conclusion $\nu_j\theta_j(x)$.
A nonleaf has the modus-ponens tag and two ordered children.
If their conclusions are $\alpha$ and
$\alpha\Rightarrow_j\beta$, its conclusion is $\beta$.

Let $\mathrm{Prf}_j$ be the set of codes of these trees.
Tree shape, formula syntax, and premise matching can be checked
by recursive operations on the numerical codes. Membership in
$A_j$ and $B_j$ is semidecidable. Dovetailing the required semidecision
procedures over all tree codes enumerates exactly $\mathrm{Prf}_j$.
Hence $\mathrm{Prf}_j$ is a c.e.\ carrier.
Reading conclusions from leaves to root gives a morphism
$c_j:\mathrm{Prf}_j\to S_j$ of $\mathcal C_{\mathrm{rec}}$.

For $i=1,2$, replace formula parameters by their images under
$\sigma_i$, positive or negative labels $x$ by $r_i(x)$, and
retain every node tag. The resulting map
$\rho_i:\mathrm{Prf}_0\to\mathrm{Prf}_i$ is a morphism of
$\mathcal C_{\mathrm{rec}}$.
The inclusions $r_i(A_0)\subseteq A_i$ and $r_i(B_0)\subseteq B_i$
ensure that relabeling takes positive and negative axiom leaves to
valid axiom leaves. The map $\sigma_i$ commutes with the formula
constructors and hence preserves the equalities required at
modus-ponens nodes. Thus $\rho_i$ preserves validity.
Induction on proofs gives
\[
 c_i\rho_i=\sigma_i c_0.
\]
The sets $\mathrm{Prf}_j$ and the maps $\rho_i$ define the functor
$\mathrm{Prf}:J\to\mathcal C_{\mathrm{rec}}$ by
$\mathrm{Prf}(j)=\mathrm{Prf}_j$ and
$\mathrm{Prf}(0\to i)=\rho_i$. The maps $c_j$ form a
natural transformation $c:\mathrm{Prf}\to S$, hence a morphism of
$\mathcal E$.

\begin{theorem}\label{thm:fork-proof-presentation}
For disjoint restriction idempotents $a,b$ on $V$ in $\mathcal A$, the syntax and
proof objects above define a restriction idempotent $t_0=\ran(c)$ on $S$ in
$\mathcal A$. The subfunctor $D_{t_0}\subseteq S$ is the least c.e.\
subfunctor containing all instances of the three logical schemes and the
positive and negative atomic axioms, and closed under modus ponens.
Moreover,
\[
 a\le\theta^*t_0,\qquad
 b\le\theta^*\nu^*t_0,\qquad
 t_0\,\nu^*t_0=0_S.
\]
At each component, these conclusions form a consistent theory of
the propositional calculus specified above.
\end{theorem}
\begin{proof}
Let $T_j=c_j(\mathrm{Prf}_j)\subseteq S_j$.
These sets are c.e.\ images of c.e.\ carriers under
morphisms of $\mathcal C_{\mathrm{rec}}$. The conclusion equation gives
$\sigma_i(T_0)\subseteq T_i$, so they form a c.e.\
subfunctor $T\subseteq S$. Its partial identity is
$\ran(c)$: it receives $c$, and the subfunctor of every restriction
idempotent receiving $c$ contains its component images. Thus $T=D_{t_0}$.

Each leaf supplies a proof of its axiom instance. Given proofs
of $\alpha$ and $\alpha\Rightarrow_j\beta$, adjoining a
modus-ponens root supplies a proof of $\beta$. Thus $T$ is
closed under modus ponens. If another c.e.\ subfunctor contains all
the axiom instances and is closed under modus ponens, induction
on each proof puts its conclusion in that subfunctor. This
proves leastness. The one-node positive and negative proofs define
morphisms $A\to\mathrm{Prf}$ and $B\to\mathrm{Prf}$ of $\mathcal E$.
Their composites with $c$ are $\theta|_A$ and
$\nu\theta|_B$, respectively, proving the first two inequalities.

For consistency, fix a component $j$ and suppose two proofs
$p,q\in\mathrm{Prf}_j$ conclude $\varphi$ and $\nu_j\varphi$.
Let $L\subseteq A_j$ be the finite set of atom labels at
positive leaves of either proof. Assign truth value $1$ to an
atom precisely when its label belongs to $L$, and assign $0$
to every other atom. Extend these values to formulas using
Boolean negation and implication.
All positive leaves of both proofs are true. Each negative
leaf has a label in $B_j$, which is disjoint from $A_j$ and
hence from $L$, so that leaf is true as well.

The three logical schemes are Boolean tautologies, and modus
ponens preserves truth. Induction on
$p$ and $q$ would therefore make $\varphi$ and its negation
both true, a contradiction.

Consequently $T_j\cap\nu_j^{-1}(T_j)=\varnothing$ for every
$j$. Restriction idempotents, inverse images, and their meets are componentwise
in $\mathcal A$, so this is precisely $t_0\nu^*t_0=0_S$.
\end{proof}

\subsection{The universal property of the proof object}\label{app:proof-initial}

The conclusion morphism $c:\mathrm{Prf}\to S$ is an object of the
slice category $\mathcal E/S$. We characterize it by the operations
that form axiom leaves and apply modus ponens.

The logical schemes define morphisms
$H_1,H_3:S\times S\to S$ and $H_2:S^3\to S$ of $\mathcal E$.
Let
\[
 \mathrm{Ax}=(S\times S)+S^3+(S\times S)+A+B,
 \qquad
 \lambda=[H_1,H_2,H_3,\theta|_A,\nu\theta|_B]:
       \mathrm{Ax}\longrightarrow S.
\]
An element of $\mathrm{Ax}_j$ records the tag and parameters of an
axiom instance; $\lambda_j$ gives its conclusion.
Write $\pi_1,\pi_2:S\times S\to S$ for the projections and put
\[
 \mu=\langle\pi_1,\Rightarrow\rangle:S\times S\longrightarrow S\times S.
\]
Thus $\mu_j(\alpha,\beta)=(\alpha,\alpha\Rightarrow_j\beta)$.

For an object $q:Q\to S$ of $\mathcal E/S$, form the pullback
\[
\begin{tikzcd}[column sep=large,row sep=large]
 M(Q,q) \arrow[r,"v_q"] \arrow[d,"w_q"'] &
 Q\times Q \arrow[d,"q\times q"] \\
 S\times S \arrow[r,"\mu"'] & S\times S.
\end{tikzcd}
\]
At component $j$, its elements are tuples
$(p_1,p_2,\alpha,\beta)$ satisfying
$q_j(p_1)=\alpha$ and
$q_j(p_2)=\alpha\Rightarrow_j\beta$.
The morphism
$\kappa_q=\pi_2w_q:M(Q,q)\to S$ gives the conclusion $\beta$.

Define an endofunctor $\mathsf H:\mathcal E/S\to\mathcal E/S$ by
\[
 \mathsf H(Q,q)=
 \bigl(\mathrm{Ax}+M(Q,q),[\lambda,\kappa_q]\bigr).
\]
For a slice morphism $h:(Q,q)\to(R,r)$, so that $rh=q$,
the pullback property gives a unique
$M(h):M(Q,q)\to M(R,r)$ satisfying
\[
 v_rM(h)=(h\times h)v_q,\qquad w_rM(h)=w_q.
\]
Put $\mathsf H(h)=1_{\mathrm{Ax}}+M(h)$.
The second equation gives $\kappa_rM(h)=\kappa_q$, so this is a
slice morphism. Uniqueness in the pullback proves the identity and
composition laws for $\mathsf H$.

Forming an axiom leaf defines $\ell:\mathrm{Ax}\to\mathrm{Prf}$.
Adjoining a modus-ponens root to two proofs defines
$m:M(\mathrm{Prf},c)\to\mathrm{Prf}$.
The implication constructor is injective at each component, so the two
children of a valid modus-ponens node determine its consequent uniquely.
The maps $\ell$ and $m$ are morphisms of $\mathcal E$: the operations on tree codes
are recursive, and relabeling commutes with their constructors.
They satisfy $c\ell=\lambda$ and $cm=\kappa_c$.
Consequently
\[
 \delta=[\ell,m]:\mathsf H(\mathrm{Prf},c)
                     \longrightarrow(\mathrm{Prf},c)
\]
is an $\mathsf H$-algebra in $\mathcal E/S$.

\begin{proposition}\label{prop:fork-proof-initial}
The algebra $\delta$ is initial. For every
$\mathsf H$-algebra
$\eta:\mathsf H(Q,q)\to(Q,q)$ in $\mathcal E/S$, there is a unique
morphism $u:\mathrm{Prf}\to Q$ in $\mathcal E$ such that
\[
 qu=c,\qquad
 u\delta=\eta\bigl(1_{\mathrm{Ax}}+M(u)\bigr).
\]
Equivalently, the following square commutes in $\mathcal E/S$:
\[
\begin{tikzcd}[column sep=large,row sep=large]
 \mathsf H(\mathrm{Prf},c) \arrow[r,"\delta"]
       \arrow[d,"\mathsf H(u)"'] &
 (\mathrm{Prf},c) \arrow[d,"u"] \\
 \mathsf H(Q,q) \arrow[r,"\eta"'] & (Q,q).
\end{tikzcd}
\]
\end{proposition}
\begin{proof}
At component $j$, evaluate a finite proof tree from its leaves to
its root. At an axiom leaf with datum $\xi\in\mathrm{Ax}_j$, apply
$\eta_j$ to the corresponding element of the first summand.
Since $\eta$ is a slice morphism, the result lies in
$q_j^{-1}(\{\lambda_j(\xi)\})$.

At a modus-ponens node, suppose the two subtrees conclude
$\alpha$ and $\alpha\Rightarrow_j\beta$.
By induction, their values $p_1,p_2\in Q_j$ satisfy
$q_j(p_1)=\alpha$ and $q_j(p_2)=\alpha\Rightarrow_j\beta$. Hence
$(p_1,p_2,\alpha,\beta)$ belongs to $M(Q,q)_j$.
Apply $\eta_j$ to its image in the second summand. The result
lies in $q_j^{-1}(\{\beta\})$.

This recursion defines $u_j:\mathrm{Prf}_j\to Q_j$ and proves
$q_ju_j=c_j$. Each step uses a partial recursive function total
on its specified carrier. The tree is finite, so its evaluation
is total on $\mathrm{Prf}_j$ and is a morphism of
$\mathcal C_{\mathrm{rec}}$.
Naturality of $\eta$, $\lambda$, and the pullback projections
shows inductively that evaluation commutes with relabeling.
Thus the maps $u_j$ form a morphism $u:\mathrm{Prf}\to Q$ of
$\mathcal E$. Their defining recursion is the algebra equation.
Any other algebra morphism agrees with $u$ on leaves and then
at each modus-ponens node by induction on the tree, proving uniqueness.
\end{proof}

For $\varphi\in S_j$, the fiber $c_j^{-1}(\{\varphi\})$ consists
of the derivation trees with conclusion $\varphi$. The range
$\ran(c)$ is the partial identity on the subfunctor of derivable formulas.


\subsection{First-order derivations and sentence proofs}
\label{app:first-order-proofs}

Fix a recursive first-order language and a recursive coding of its
formulas. Let $T$ be a theory in this language whose nonlogical axioms
are sentences. Let
$\mathcal D$ be $\mathcal C_{\mathrm{rec}}$ for a c.e.\ axiom set,
or $\operatorname{Tot}(\operatorname{Split}(\mathcal E_a))$ for an
$\mathcal E_a$-semidecidable axiom set. In the latter case we use the
\hyperref[def:miller-splitting]{carrier presentation of the splitting}.
Use a sound and complete
Hilbert calculus with a recursive set of logical axioms and the rules
of modus ponens and generalization.
Write $\mathsf F$ for the recursive carrier of all formulas,
$\mathsf S\hookrightarrow\mathsf F$ for its sentence subobject, and
$\mathsf L,\mathsf T\hookrightarrow\mathsf F$ for the logical and
nonlogical axiom carriers.

The objects of rule instances are
\[
 R_0=\mathsf L+\mathsf T,\qquad
 R_1=\N\times\mathsf F,\qquad R_2=\mathsf F\times\mathsf F.
\]
For $k=0,1,2$, let $p_k:R_k\to\mathsf F^k$ and
$q_k:R_k\to\mathsf F$ give the premises and conclusion. Here
$\mathsf F^0=1$, $p_0$ is the terminal map, and $q_0$ is the copairing
of the two axiom inclusions. For generalization and modus ponens, put
\[
\begin{aligned}
 p_1(v,\varphi)&=\varphi,&
 q_1(v,\varphi)&=\forall v\,\varphi,\\
 p_2(\varphi,\psi)&=(\varphi,\varphi\Rightarrow\psi),&
 q_2(\varphi,\psi)&=\psi.
\end{aligned}
\]
Variables are indexed by $v\in\N$. Generalization has no additional
side condition because all nonlogical assumptions are sentences.
The displayed connectives are constructors on formula codes.
These maps are morphisms of $\mathcal D$.

The products, pullbacks, and finite coproducts below exist in
$\mathcal D$. In the chosen carrier presentations they are formed
by finite products, equality tests, and tagged unions of the specified
semidecidable sets.

For $k=0,1,2$ and $(Q,q)\in\mathcal D/\mathsf F$, form the pullback
of $p_k$ and $q^k$. Write
\[
\begin{gathered}
 B_k(Q,q)=R_k\times_{\mathsf F^k}Q^k,
 \\
 a_{k,q}:B_k(Q,q)\to R_k,\qquad
 b_{k,q}:B_k(Q,q)\to Q^k
\end{gathered}
\]
for its object and projections, so that
$p_ka_{k,q}=q^kb_{k,q}$.
Define $\mathsf H_T:\mathcal D/\mathsf F\to\mathcal D/\mathsf F$
on objects by
\[
 \mathsf H_T(Q,q)=
 \left(\coprod_{k=0}^2 B_k(Q,q),
       [q_ka_{k,q}]_{k=0}^2\right).
\]
For a slice morphism $h:(Q,q)\to(Q',q')$, the identity $q'h=q$
and the pullback property give a unique
$B_k(h):B_k(Q,q)\to B_k(Q',q')$ satisfying
\[
 a_{k,q'}B_k(h)=a_{k,q},\qquad
 b_{k,q'}B_k(h)=h^kb_{k,q}.
\]
Put $\mathsf H_T(h)=\coprod_{k=0}^2 B_k(h)$.
The first equation makes this a slice morphism. Pullback uniqueness
gives the identity and composition laws.

\begin{proposition}\label{prop:first-order-proof-initial}
Let $\mathsf{Der}_T$ be the carrier of codes of finite rooted ordered
trees whose nodes are labeled by rule instances. A node labeled by
$r\in R_k$ has exactly $k$ children, whose ordered tuple of conclusions
is $p_k(r)$; its conclusion is $q_k(r)$. Let
$d_T:\mathsf{Der}_T\to\mathsf F$ give their root conclusions.
Adjoining a rule at the root defines an algebra
\[
 \delta_T:\mathsf H_T(\mathsf{Der}_T,d_T)
       \longrightarrow(\mathsf{Der}_T,d_T).
\]
This algebra is initial among $\mathsf H_T$-algebras in
$\mathcal D/\mathsf F$.
The proof object with sentence conclusions is the pullback
\[
\begin{tikzcd}[column sep=large,row sep=large]
 \mathrm{Prf}_T \arrow[r] \arrow[d,"c_T"'] &
 \mathsf{Der}_T \arrow[d,"d_T"] \\
 \mathsf S \arrow[r,hook] & \mathsf F.
\end{tikzcd}
\]
Viewed in $\operatorname{Par}(\mathcal C_{\mathrm{rec}},\operatorname{Mono})$
or $\operatorname{Split}(\mathcal E_a)$, respectively, the map $c_T$
has range equal to the partial identity on the theorem set of $T$.
\end{proposition}
\begin{proof}
Tree shape, syntactic well-formedness of the parameters, and premise
matching are recursively decidable. Membership in the nonlogical
axiom carrier is semidecidable in the stated sense. Each tree has
only finitely many nonlogical axiom labels, so the valid tree codes
form an object of $\mathcal D$.
The conclusion map and the operation adjoining a rule at the root are
recursive on their specified carriers. In the splitting we restrict them by the partial
identities on those carriers. The resulting total morphisms satisfy the
conclusion equations.

For an algebra $\eta:\mathsf H_T(Q,q)\to(Q,q)$, evaluation from
leaves to root defines $u:\mathsf{Der}_T\to Q$.
Induction on the tree proves termination, $qu=d_T$, and the
algebra-morphism equation. The same induction proves uniqueness. In
$\operatorname{Tot}(\operatorname{Split}(\mathcal E_a))$, first run the
semidecision procedure for membership in the carrier of valid trees.
Evaluation then has exactly that domain; its value is independent of
the binary representation because each algebra operation has this property. Thus the unique evaluation map belongs to
$\mathcal D$ in both cases.

The pullback selects precisely the trees whose root conclusions are
sentences. Closedness is recursive on formula codes, so this selection
preserves semidecidability. Its image under the conclusion map is the
set of theorems.
\end{proof}

